\section*{INTRODUCTION}

In the beginning, when Dirac, Jordan, Heisenberg, and Pauli created the quantum theory of fields, it was not expected that it would provide a consistent description of Nature. After all, it was only a quantized version of the classical theory of Maxwell and Lorentz, a theory which was well known to be afflicted with diseases arising from the infinite electromagnetic inertia of point particles. Many physicists were of the opinion that any project to make the theory's mathematical foundation more rigorous was probably ill-advised; first the classical foundation should be set right. Such alterations might so change the basis of the theory that a mathematically rigorous discussion of any preceding version would be entirely irrelevant. More recently, it has been suggested that the trouble is that the theory is too modest; it is not designed to predict the masses of the elementary particles or the values of the coupling constants, and should be fundamentally changed with this in view.

However, attempts to go beyond the theory foundered again and again. What successes were achieved were either phenomenological, or were due to systematic developments of the original formalism. But the quantum theory of fields never reached a stage where one could say with confidence that it was free from internal contradictions---nor the converse. In fact, the Main Problem of quantum field theory turned out to be to kill it or cure it: either to show that the idealizations involved in the fundamental notions of the theory (relativistic invariance, quantum mechanics, local fields, etc.) are incompatible in some physical sense, or to recast the theory in such a form that it provides a practical language for the description of elementary particle dynamics.

The last ten years have seen a number of attempts to meet the situation head on. (The physicists who have engaged in this kind of work are sometimes dubbed the Feldverein. Cynical observers have compared them to the Shakers, a religious sect of New England who built solid barns and led celibate lives, a non-scientific equivalent of proving rigorous theorems and calculating no cross sections.) These efforts have not yet led to a solution of the Main Problem, but they have yielded a number of by-products, very general insights into the structure of a field theory. The present book is devoted to an exposition of some of these general results, the physical ideas they embody, and the mathematics necessary for their proofs.

We have included only results which have a certain definitive character. In particular, this has resulted in the omission of a description of attempts to establish a connection with the important work of Lehmann, Symanzik, and Zimmermann and others on time-ordered and retarded functions and their connection with collision theory. A great deal more work will be necessary before these subjects can be properly understood and put on a rigorous basis. Although the connection with Lehmann, Symanzik, and Zimmermann is not yet firmly established, a rigorous collision theory (based on the axioms of Chapter 3 of the present book) has been set up by D. Ruelle along lines laid down by R. Haag. The omission of this theory from this book will be mitigated by the availability of the excellent book by R. Jost, which will contain a full account (\emph{The General Theory of Quantized Fields}, American Mathematical Society, 1965).

The first chapter contains a summary of the transformation properties of physical states in relativistic quantum mechanics. It is assumed that the reader has had an introduction to Hilbert space and its application to the description of states in quantum mechanics. It is probably worthwhile to point out to younger readers that the simplicity of the transformation laws of the physical vacuum and one-particle states under Lorentz transformations, well known today, was buried in the quantum field theory of fifteen years ago under a kitchen midden of difficult and ambiguous formalism. The task of the first chapter is to provide a language in which physical states with simple transformation properties have a simple description. For example, the concepts of bare mass and bare vacuum need not be and are not introduced.

The second chapter is an exposition of the mathematical tools used in the following. Technical details of some proofs have been omitted but an attempt has been made to get across the main mathematical ideas. The theorems are stated precisely. The presentation presupposes nothing that an undergraduate physics major has not met.

The third chapter defines the notion of field as used in this book. It is shown that a field theory is defined by the vacuum expectation values of products of field operators. While this chapter is essentially self-contained, a brush with elementary quantum field theory at the level of, say, Part II of Schweber's book
\begingroup
\renewcommand{\thefootnote}{\ensuremath{\dagger}}%
\footnote{S. Schweber, \emph{An Introduction to Relativistic Quantum Field Theory}, Harper and Row, New York, 1961.}
\endgroup
might be a help.

In Chapter 4 the three preparatory chapters are applied to get some general theorems of quantum field theory, of which the \emph{PCT} theorem and the theorem on the connection of spin with statistics are the best known.

The reader who wants to get on with the systematic discussion of quantum field theory could well begin with Chapter 3 and only go back to Chapters 1 and 2 when he finds it necessary to fill in details.

Each chapter has been equipped with a bibliography to guide the reader to relevant literature. No attempt at completeness has been made. The notation is standard: Theorem 3-1 refers to the first theorem of the third chapter, and similarly for equations. Halmos notation \(\blacksquare\) has been used to signify the end of a proof.
